\documentclass{article}
\usepackage[utf8]{inputenc}
\usepackage[T1]{fontenc}

\title{Sprawozdanie z projektu na temat wskaźnika giełdowego MACD}
\author{Piotr Kaczorowski 197736}
\date{\today}

\begin{document}

\maketitle

\section{Wstęp teoretyczny}

Wskaźnik MACD został opracowany przez Geralda Appela w latach\
70-tych XX wieku. W celu ukazywania zmian w trendach cenowych akcji\
na giełdzie. Może być wykorzystywany również przy analizie innych\
instrumentów finansowych, takich jak kryptowaluty czy surowce.\
Sam MACD jest różnicą dwóch średnich kroczących (EMA - \
Exponential Moving Average), zwykle\
będących średnimi z 12 (EMA$_{12}$) i 26 ostatnich okresów (EMA$_{26}$). \
W celu wyznaczenia sygnałów kupna i sprzedaży stosuje się\
linię sygnałową, będącą średnią kroczącą (zwykle z 9 okresów) MACD.\
Przecięcia linii MACD z liniami sygnałowymi na wykresie są interpretowane\
jako sygnały kupna i sprzedaży.

Celem niniejszego projektu jest implementacja wskaźnika MACD w języku\
Python, a następnie przetestownanie go na na przykładowych danych\
giełdowych oraz zaprezentowanie wyników w formie wykresów jak również\
opracowania wniosków na ich podstawie. Podczas implementowania wskaźnika\
MACD została wykorzystana biblioteka Pandas do wczytania odpowiednich danych\
z pliku CSV zawierającego dane giełdowe oraz biblioteka Matplotlib do\
wygenerowania wykresów na podstawie otrzymanych wyników i danych wejściowych.\
Kod źródłowy projektu jest zapisany w pliku \emph{MACD.py} (a do debugowania\
i testowania kodu został wykorzystany plik \emph{debug MACD.ipynb})


\section{Dane przykładowe i implementacja}

Dane przykładowe zostały pobrane z serwisu stooq.pl w formie pliku CSV.\
Zawierają one dane giełdowe na temat wartości WIG 20 w okresie od 16.04.2014\
do 16.10.2018. Dane te zostały wczytane do programu za pomocą biblioteki\
Pandas, a następnie zostały wydzielone interesujące nas dane,\
które są cenami zamknięcia akcji. Dane te zostały przekazane do funkcji\
implementującej wyliczanie wskaźnika EMA. Funkcja ta implementuje\
wyliczanie średniej kroczącej wykorzystując wzór:
\begin{equation}
    EMA_n = \frac{p_0+(1-\alpha)*p_1+(1-\alpha)^2*p_2+...+(1-\alpha)^n*p_n}{1+(1-\alpha)+(1-\alpha)^2+...+(1-\alpha)^n}
\end{equation}

p$_{0}$,p$_{1}$ ... p$_{n}$ - kolejne wartości z których liczona jest średnia

${\alpha}$ = ${\frac{2}{n+1}}$%, jest to wyliczane w funkcji \emph{alpha_calc}      

${EMA_n}$ - średnia krocząca z n elementów, im starszy tym ma mniejsze znaczenie
\newline
\newline
Z implementacji powyższego wzoru (1) wyliczane są tablice zawierające ${EMA_{12}}$\
i ${EMA_{26}}$. Następnie jest obliczany wskaźnik MACD z obu powyższych wyników odejmując od siebie\
odpowiadające wartości Po obliczeniu wskaźnika MACD zostaje wyliczenie wartości linii\
SIGNAL, obliczona ona jest za pomocą funkcji implementującej ${EMA_n}$ wywołanej z\
innymi parametrami - jako zbiór danych zostaje przesłany zbiór wartości MACD\
a \emph{n} w tym wywołaniu jest ustawiane na 9.
\newline
Kolejnym krokiem jest wyznaczenie punktów kupna i sprzedaży, to zadanie wykonuje\
funkcja \emph{intersectionPoints} zwracająca kolejno: listę indeksów punktów przecięcia\
listę tupli kupna zawierającą indeks oraz wartość MACD w tym indeksie\
(różnica pomiędzy MACD i SIGNAL jest pomijalna albo niemożliwa\
do usunięcia w implementacji - bo oba wskaźniki muszą mieć wartość identyczną, jednak\
w rzeczywistych warunkach zakresu wartości niekoniecznie jest to możliwe oraz\
przecięcia mogą występować pomiędzy dwoma polami wyrażonymi w MACD i SIGNAL),\
ostatnia jest lista sprzedaży, mająca identyczną strukturę do listy kupna.


\section{Rezultat na przykładowych danych}


\section{Conclusion}
Conclusion of the document.

\end{document}